\documentclass[10pt,a4paper]{amsart}
\usepackage{amsmath,amssymb,amsthm,graphicx}
\usepackage{verbatim}
\newenvironment{reference}{\comment}{\endcomment}
%\newenvironment{reference}{}{}
\newenvironment{slogan}{\comment}{\endcomment}
\newenvironment{history}{\comment}{\endcomment}

% For commutative diagrams we use Xy-pic
\usepackage[all]{xy}

% We use 2cell for 2-commutative diagrams.
\xyoption{2cell}
% Theorem environments.
%
\theoremstyle{plain}
\newtheorem{theorem}[subsection]{Theorem}
\newtheorem{proposition}[subsection]{Proposition}
\newtheorem{lemma}[subsection]{Lemma}

\theoremstyle{definition}
\newtheorem{definition}[subsection]{Definition}
\newtheorem{example}[subsection]{Example}
\newtheorem{exercise}[subsection]{Exercise}
\newtheorem{situation}[subsection]{Situation}

\theoremstyle{remark}
\newtheorem{remark}[subsection]{Remark}
\newtheorem{remarks}[subsection]{Remarks}

\numberwithin{equation}{subsection}

% Macros
%
\def\lim{\mathop{\mathrm{lim}}\nolimits}
\def\colim{\mathop{\mathrm{colim}}\nolimits}
\def\Spec{\mathop{\mathrm{Spec}}}
\def\Hom{\mathop{\mathrm{Hom}}\nolimits}
\def\Ext{\mathop{\mathrm{Ext}}\nolimits}
\def\SheafHom{\mathop{\mathcal{H}\!\mathit{om}}\nolimits}
\def\SheafExt{\mathop{\mathcal{E}\!\mathit{xt}}\nolimits}
\def\Sch{\mathit{Sch}}
\def\Mor{\mathop{\mathrm{Mor}}\nolimits}
\def\Ob{\mathop{\mathrm{Ob}}\nolimits}
\def\Sh{\mathop{\mathit{Sh}}\nolimits}
\def\NL{\mathop{N\!L}\nolimits}
\def\CH{\mathop{\mathrm{CH}}\nolimits}
\def\proetale{{pro\text{-}\acute{e}tale}}
\def\etale{{\acute{e}tale}}
\def\QCoh{\mathit{QCoh}}
\def\Ker{\mathop{\mathrm{Ker}}}
\def\Im{\mathop{\mathrm{Im}}}
\def\Coker{\mathop{\mathrm{Coker}}}
\def\Coim{\mathop{\mathrm{Coim}}}

% Boxtimes
%
\DeclareMathSymbol{\boxtimes}{\mathbin}{AMSa}{"02}

%
% Macros for moduli stacks/spaces
%
\def\QCohstack{\mathcal{QC}\!\mathit{oh}}
\def\Cohstack{\mathcal{C}\!\mathit{oh}}
\def\Spacesstack{\mathcal{S}\!\mathit{paces}}
\def\Quotfunctor{\mathrm{Quot}}
\def\Hilbfunctor{\mathrm{Hilb}}
\def\Curvesstack{\mathcal{C}\!\mathit{urves}}
\def\Polarizedstack{\mathcal{P}\!\mathit{olarized}}
\def\Complexesstack{\mathcal{C}\!\mathit{omplexes}}
% \Pic is the operator that assigns to X its picard group, usage \Pic(X)
% \Picardstack_{X/B} denotes the Picard stack of X over B
% \Picardfunctor_{X/B} denotes the Picard functor of X over B
\def\Pic{\mathop{\mathrm{Pic}}\nolimits}
\def\Picardstack{\mathcal{P}\!\mathit{ic}}
\def\Picardfunctor{\mathrm{Pic}}
\def\Deformationcategory{\mathcal{D}\!\mathit{ef}}
\usepackage{cleveref}
\crefname{originalreference}{original reference}{original references}
\Crefname{originalreference}{Original reference}{Original references}
\makeatletter
\newlabel{lemma-get-triangle}{{09P6}{source}{}{sourcepointer-1}{}}
\newlabel{lemma-get-triangle@cref}{{[originalreference][0][]09P6}{[1][0][]source}{}{}{}}
\makeatother
\title{A functor between differential graded categories with axioms (A), (B), (C) induces an exact functor}
\author{The Stacks Project Authors}
\date{stacks-project:dga.tex:local-snapshot-20261012; GNU FDL 1.2 or later}
\begin{document}
\maketitle
\noindent\textit{Editorial note.} All delivered mathematical text is the author's own, copied from the pinned Stacks Project chapter file sources/stacks-dga/pinned-original-source.tex; no statement or proof step was rewritten, reordered or removed.\par\medskip
\noindent\textit{Editorial note.} The counted proof is the author's complete argument: the preservation of admissible short exact sequences (proof line 4716-4718) is the standard additivity argument: a functor of differential graded categories is R-linear, hence additive, hence preserves graded biproducts and commutes with differentials, so termwise-split sequences map to termwise-split sequences; the boundary agreement is computed explicitly in the span; the boundary construction is Lemma lemma-get-triangle (in-source, navigation stub).\par\medskip
\noindent\textit{Editorial note.} The label lemma-functor-between-ABC names the counted statement, and the proof environment immediately following it in the source is delivered verbatim as its proof.\par\medskip
\noindent\textit{Editorial note.} The counted statement is carried by the author's own lemma environment in the chapter source.\par\medskip
\noindent\textit{Allowed standard prerequisites (editorial).}\begin{itemize}
\item The ABC axioms and the boundary construction. Same unit, Situation situation-ABC and Lemma lemma-get-triangle. The proof uses the boundary morphism delta = pi d(s) attached to an admissible short exact sequence by the construction of Lemma lemma-get-triangle, and the axioms (A), (B), (C) enjoyed by both categories.
\item Shift compatibility hypothesis. Hypothesis of the lemma: F(x[1]) = F(x)[1]. The hypothesis identifies the target of F(delta) with F(x)[1], so the image of the boundary morphism lands in the right place.
\end{itemize}\medskip
\ifdefined\hypertarget\hypertarget{sourcepointer-1}{}\fi
\noindent\textit{Source reference.} Original reference 09P6: sources/stacks-dga/pinned-original-source.tex, source lines 3708--3708. Prerequisite pointer invoked at proof line(s) 4724; it names a standard fact of the same chapter that the counted argument uses, and every step of the counted proof built on it is retained verbatim in the delivered body.\par\medskip
\begin{lemma}
\label{lemma-functor-between-ABC}
Let $R$ be a ring. Let $F : \mathcal{A} \to \mathcal{B}$ be a functor
between differential graded categories over $R$ satisfying axioms
(A), (B), and (C) such that $F(x[1]) = F(x)[1]$.
Then $F$ induces an exact functor
$K(\mathcal{A}) \to K(\mathcal{B})$ of triangulated categories.
\end{lemma}

\begin{proof}
Namely, if $x \to y \to z$ is an admissible short exact sequence
in $\text{Comp}(\mathcal{A})$, then $F(x) \to F(y) \to F(z)$
is an admissible short exact sequence in $\text{Comp}(\mathcal{B})$.
Moreover, the ``boundary'' morphism $\delta = \pi\text{d}(s) : z \to x[1]$
constructed in Lemma \ref{lemma-get-triangle} produces the morphism
$F(\delta) : F(z) \to F(x[1]) = F(x)[1]$ which is equal to the boundary
map $F(\pi) \text{d}(F(s))$ for the admissible short exact sequence
$F(x) \to F(y) \to F(z)$.
\end{proof}

\end{document}
